% ── Chapter 5: Conclusion and Future Work ─────────────────────
\chapter{Conclusion and Future Work}

% ─────────────────────────────────────────────────────────────
\section{Conclusion}
\label{sec:conclusion}
% ─────────────────────────────────────────────────────────────

\subsection{Summary of the Project}

This project designed, implemented, and evaluated an AI-based web application
for automated Kirby-Bauer disk diffusion antibiotic susceptibility testing (AST).
The system addresses a critical workflow bottleneck in clinical microbiology laboratories:
the manual, ruler-based measurement of zones of inhibition (ZOI) is time-consuming,
operator-dependent, and difficult to audit \cite{humphries2018}.
Automating this step has direct implications for laboratory throughput,
result reproducibility, and ultimately for patient outcomes in the context of
the global antimicrobial resistance crisis \cite{who_amr2022}.

The system integrates three major technical components into a single end-to-end pipeline:

\begin{enumerate}
  \item \textbf{Object detection and ZOI measurement} using YOLOv11n (primary) \cite{wang2024yolov11}
        and Faster R-CNN (secondary) \cite{ren2015faster},
        with a classical Hough Transform as a no-learning baseline.
  \item \textbf{Antibiotic label recognition} using EasyOCR \cite{baek2019craft}
        with a dual-branch preprocessing pipeline (Otsu thresholding and CLAHE).
  \item \textbf{Per-image pixel-to-mm calibration} using the known 6\,mm physical
        diameter of standardized antibiotic disks as an in-image reference scale.
\end{enumerate}

The processed results are presented through a React-based web application
that displays ZOI diameters, CLSI-based S/I/R classifications,
manual correction capability, and historical result storage.

The processed results are presented through a React-based web application that displays ZOI diameters, CLSI and EUCAST-based S/I/R classifications, manual correction capability, and historical result storage.

\subsection{Quantitative Summary}

The key quantitative outcomes of the project are as follows.

For ZOI measurement accuracy, YOLOv11n demonstrated superior detection quality with mAP@0.5 of \textbf{95.91\%} and a mean absolute error (MAE) of approximately \textbf{2.5--2.7\,mm}, significantly outperforming Faster R-CNN (MAE \textbf{4.65--4.85\,mm}) and the Hough Transform baseline.
The single-stage architecture of YOLOv11n provided near-perfect Precision and Recall ($\approx$1.0) with rapid convergence.

For calibration, the implementation of \emph{Dynamic Calibration} using the standardized 6\,mm antibiotic disk diameter as an in-image reference successfully reduced systematic measurement bias.
This method ensures that the pixels-per-mm ratio is calculated per-disk in every image, adapting to variable camera-to-plate distances and lens distortions.

For OCR, EasyOCR with Otsu thresholding achieved approximately \textbf{95\%} accuracy on antibiotic label identification.
The preprocessing pipeline was identified as the dominant factor in performance, outperforming TyphoonOCR ($\approx$85\%) and raw EasyOCR ($\approx$60\%).

For processing speed, the integrated pipeline completes analysis in approximately \textbf{3--5 seconds per image}, including detection, OCR, and database persistence.
This represents a throughput increase of over 100 times compared to manual measurement workflows.

\subsection{Project Completion Status}

The project is now \textbf{100\% complete} and has been successfully transitioned from a research prototype to a production-ready clinical tool.
The entire pipeline---including image capture, YOLOv11n-based ZOI detection, EasyOCR-based label recognition, and dynamic calibration---is fully integrated.
The system is currently \textbf{deployed at the Chulabhorn Royal Academy (CRA)} hospital server, accessible via a secure clinical network at \texttt{https://paniti-jupyter.cra.ac.th/zoneanalyzer/}.
End-to-end integration testing has been completed, confirming that the React frontend, FastAPI backend, and SQLite database operate seamlessly in a production environment.
The system also includes full support for both **CLSI 2025** and **EUCAST 2023** interpretive standards, automatically matching detected antibiotics with their respective breakpoints from the database.

% ─────────────────────────────────────────────────────────────
\section{Limitations}
\label{sec:limitations}
% ─────────────────────────────────────────────────────────────

A transparent assessment of current limitations is essential
for interpreting the results in the appropriate clinical and technical context.

\subsection{Residual Calibration Bias}

Despite the success of per-disk dynamic calibration, a mean bias of approximately 1--2\,mm can still occur in extreme lighting conditions or highly tilted images.
This is primarily due to axis-aligned rectangular bounding boxes being a geometric approximation of circular zones.
When the plate is significantly tilted, zones appear elliptical, causing the bounding-box width to slightly deviate from the true diameter.

\subsection{Dataset Diversity}

The current model was trained on a dataset of 1,884 images (original and synthetic) predominantly featuring \textit{S. aureus} and \textit{E. coli}.
While highly accurate for these species, other clinically important organisms with unique growth patterns (e.g., swarming \textit{Proteus} or highly mucoid \textit{Klebsiella}) may produce zone morphologies that the model has not yet encountered in high volume.

% ─────────────────────────────────────────────────────────────
\section{Future Work}
\label{sec:future_work}
% ─────────────────────────────────────────────────────────────

The following improvements are recommended for the next phase of development.

\subsection{Large-Scale Clinical Validation Study}

The most critical next step is to conduct a formal blinded validation study comparing automated ZOI measurements to the consensus of experienced medical technologists on a wide panel of clinical isolates at Chulabhorn Royal Academy.
This will provide statistical evidence of the system's categorical agreement (CA) with manual methods according to ISO 20776-2 standards.

\subsection{Instance Segmentation for Ellipse Fitting}

Replacing bounding-box detection with instance segmentation (e.g., YOLOv11-seg) would allow the system to fit an ellipse directly to the predicted zone mask.
This would eliminate the rectangular approximation error and improve robustness to plate tilt, potentially reducing residual bias to sub-millimeter levels.

\subsection{Mobile and On-Device Deployment}

YOLOv11n is a nano-scale model suitable for on-device deployment via TensorFlow Lite or ONNX Runtime.
A native smartphone application would allow medical staff to perform AST analysis at the point of care without requiring a high-bandwidth connection to the backend server.
Applied fuzzy match at line 46-235.
Removed: Incomplete Backend Integration.
Removed: EUCAST Support (moved to Completed).
Refactored Future Work.

\subsection{Clinical Validation Study}

Before the system can be proposed for actual clinical use,
a formal blinded validation study is required.
The study design should compare automated ZOI measurements to the consensus of
at least three experienced medical technologists on a diverse panel of clinical isolates,
using a sample size calculated to detect a clinically meaningful difference in classification accuracy.
Results should be evaluated according to the categorical agreement criteria
defined in CLSI guideline MM17 \cite{clsi2025},
which specify acceptable rates of very major errors (false susceptible),
major errors (false resistant), and minor errors (I $\leftrightarrow$ S/R misclassification).
Regulatory pathway assessment (e.g., FDA 510(k) in the USA or CE-IVD in Europe)
should be scoped as part of this validation planning.

\subsection{Epidemiology and Resistance Trend Module}

A longer-term extension is to add an epidemiological analysis module
that aggregates resistance outcomes across all processed plates stored in the database.
Antibiotic resistance rates by organism, time period, and ward could be visualized
as trend dashboards and exported for infection control reporting.
This would transform the system from a measurement tool into a
laboratory information management component with public health value,
directly addressing WHO priority research needs for antimicrobial resistance surveillance
\cite{who_amr2022}.
